\section{Introduction}

Single-cell RNA-seq (scRNA-seq) has enabled detailed characterization of cellular heterogeneity within tumors, revealing that malignant cells are not transcriptionally uniform but rather exist along continuous axes of transcriptional diversity \cite{Tirosh2016}. In metastatic melanoma, the Tirosh et al. study (3ca:20111) provided a comprehensive dissection of the multicellular ecosystem, identifying distinct immune, stromal, and malignant populations \cite{Tirosh2016}. Building on this foundation, Xiao et al. (2019) demonstrated that metabolic pathways--particularly mitochondrial activity and glycolysis--are upregulated in malignant cells and that single-cell resolution captures metabolic heterogeneity invisible to bulk RNA-seq \cite{Xiao2019}. These findings have spurred interest in whether metabolic gene expression alone can define discrete transcriptional states that are biologically meaningful and reproducible across patients.

However, the claim that metabolic genes define separable cell states requires validation against multiple confounding factors. First, the Xiao2019 melanoma dataset is not an independent cohort but a differently processed subset of the same underlying Tirosh melanoma study (GSE72056) \cite{Xiao2019}. Consequently, results from Xiao2019 cannot serve as an independent validation of metabolic state separation; they instead provide a methodological benchmark against which to contrast the 3CA pipeline. Second, transcriptomic expression does not directly equate to metabolic flux. A gene's mRNA abundance reflects transcriptional output, but metabolic activity is governed by enzyme kinetics, substrate availability, and post-translational regulation. Thus, clustering based solely on metabolic gene expression may capture transcriptional programs that are not functionally equivalent to metabolic states.

To address these concerns, we ask two central questions: (1) Can metabolic genes alone separate single-cell populations into discrete, reproducible states? and (2) Do these states correspond to interpretable transcriptional metabolic programs? To answer these, we perform unsupervised clustering on 1,360 metabolic genes from 4,645 melanoma cells across 19 samples, using the 3CA data atlas as a source of the underlying single-cell RNA-seq data \cite{Gavish2023}. We evaluate clustering stability via silhouette coefficients and adjusted Rand indices across multiple seeds, and we test whether the observed separation exceeds that of matched random gene sets. Crucially, we control for library complexity--a known driver of technical variation in scRNA-seq data--by residualizing expression features and re-clustering. We also assess whether malignant cells exhibit distinct metabolic states by adjusting for both patient identity and complexity.

Finally, we evaluate the biological interpretability of any identified states through patient-blocked cross-validation and comparison with highly variable gene baselines. The patient-blocked classification tests whether a supervised classifier predicts known cell-type labels, not whether it discovers new states. K-means performs a forced partition of the data; the choice of k does not prove discrete multimodal states. Similarly, failure to establish metabolic-specific states does not imply that all observed separation is technical artifact. This study provides an exploratory reanalysis of a single cohort to evaluate the biological validity of metabolic state clustering.
